% !TEX program = xelatex
% =============================================================================
% 中文求职简历 LaTeX 模板
% Author: Kody
% Repository: https://github.com/kody1126/Chinese-resume-template-work
% Compiler: XeLaTeX
%
% 快速修改指南：
% 1. 在“个人信息头部”中替换姓名、联系方式、求职意向和头像。
% 2. 在“教育背景”中替换学校、专业、成绩排名和课程信息。
% 3. 在实习经历、项目经历、校园经历与荣誉技能模块中增删示例内容。
% 4. 如需调整样式，可修改下方的字体、页面边距与 sectioncolor 配置。
% =============================================================================
\documentclass[11pt,a4paper,fontset=none]{ctexart}

% ── 字体设置 ──────────────────────────────────────────────
% 若需在 Overleaf 使用自定义字体，可将已授权的字体文件放入 fonts/ 文件夹，
% 并按 STSong.ttf、STHeiti.ttf、STKaiti.ttf 命名；macOS 本地默认使用系统华文系列。
\IfFontExistsTF{Times New Roman}{\setmainfont{Times New Roman}}{\setmainfont{TeX Gyre Termes}}

\newif\ifuseuploadedcjkfonts
\IfFileExists{fonts/STSong.ttf}{
    \IfFileExists{fonts/STHeiti.ttf}{
        \IfFileExists{fonts/STKaiti.ttf}{\useuploadedcjkfontstrue}{}
    }{}
}{}

\newif\ifusesystemcjkfonts
\IfFontExistsTF{STSong}{
    \IfFontExistsTF{STHeiti}{
        \IfFontExistsTF{STKaiti}{\usesystemcjkfontstrue}{}
    }{}
}{}

\ifuseuploadedcjkfonts
    \setCJKmainfont[
        Path = fonts/,
        BoldFont = STHeiti.ttf,
        ItalicFont = STKaiti.ttf
    ]{STSong.ttf}
    \setCJKsansfont[Path = fonts/]{STHeiti.ttf}
    \setCJKmonofont[Path = fonts/]{STHeiti.ttf}
\else\ifusesystemcjkfonts
    \setCJKmainfont[
        BoldFont = STHeiti,
        ItalicFont = STKaiti
    ]{STSong}
    \setCJKsansfont{STHeiti}
    \setCJKmonofont{STHeiti}
\else
    \setCJKmainfont[
        BoldFont = Noto Sans CJK SC,
        ItalicFont = Noto Serif CJK SC
    ]{Noto Serif CJK SC}
    \setCJKsansfont{Noto Sans CJK SC}
    \setCJKmonofont{Noto Sans CJK SC}
\fi
\fi

% ── 页面布局 ──────────────────────────────────────────────
\usepackage{geometry}
\geometry{a4paper, left=1.5cm, right=1.5cm, top=1.5cm, bottom=1.5cm}

% ── 颜色定义 ──────────────────────────────────────────────
\usepackage{xcolor}
\definecolor{sectioncolor}{RGB}{25,78,121}
\definecolor{linkcolor}{RGB}{24,90,140}
\definecolor{textcolor}{RGB}{0,0,0}

% ── 标题格式 ──────────────────────────────────────────────
\usepackage{titlesec}
\titleformat{\section}
    {\Large\bfseries\color{sectioncolor}}
    {}
    {0em}
    {\uppercase}
    [\vspace{-0.8em}\rule{\textwidth}{1pt}]
\titlespacing{\section}{0pt}{0.1em}{0.1em}

% ── 列表格式 ──────────────────────────────────────────────
\usepackage{enumitem}
\setlist[itemize]{leftmargin=*, nosep, topsep=0.2em, itemsep=0.1em, parsep=0em}

% ── 图片与排版 ────────────────────────────────────────────
\usepackage{graphicx}
\setkeys{Gin}{draft=false}
\usepackage[export]{adjustbox}
\usepackage{calc}

% ── 表格 ─────────────────────────────────────────────────
\usepackage{tabularx}
\usepackage{multirow}

% ── 图标与超链接 ──────────────────────────────────────────
\usepackage{fontawesome5}
\usepackage{hyperref}
\hypersetup{colorlinks=true, linkcolor=linkcolor, urlcolor=linkcolor}

% ── 全局样式 ──────────────────────────────────────────────
\color{textcolor}
% 华文字体适当增加字距；Linux 上的 Noto 字体保持默认字距，避免长标题换行。
\ifuseuploadedcjkfonts
    \xeCJKsetup{CJKglue = \hskip 0.15em}
\else\ifusesystemcjkfonts
    \xeCJKsetup{CJKglue = \hskip 0.15em}
\else
    \xeCJKsetup{CJKglue = \hskip 0em}
\fi\fi
\linespread{1.00}
\pagestyle{empty}

% ── 自定义命令 ────────────────────────────────────────────
\newcommand{\daterange}[1]{{\textbf{#1}}}
\newcommand{\entrytitle}[1]{{\textbf{#1}}}

% ============================================================
\begin{document}
% ============================================================

% ── 个人信息头部 ──────────────────────────────────────────
\noindent
\begin{minipage}[t]{\textwidth}
    \vspace*{-1.8em}
    \begin{center}
        {\Huge\bfseries\color{sectioncolor} 白逸凡} \\[1.0em]
        \begin{tabular}{@{}ll@{\hspace{2.5em}}ll@{}}
            \textbf{出生日期：} & 2003.09.03 & \textbf{电话：} & 173-2791-2445 \\[0.4em]
            \textbf{个人主页：} & \href{https://blog.msuad.top/resume-msuad/}{blog.msuad.top/resume-msuad} & \textbf{邮箱：} & 1445617868@qq.com \\
        \end{tabular}
    \end{center}
\end{minipage}

\vspace{0.3em}

% ── 教育背景 ──────────────────────────────────────────────
\section*{教育背景}
\vspace{-0.3em}

\noindent
\entrytitle{中国科学院大学信息工程研究所 ｜ 网络信息与安全（硕士） ｜ 体系结构安全组（国重）}
\par\noindent\hfill \daterange{2025.09 \textasciitilde{} 2028.07（预计）}

\noindent
\entrytitle{南京邮电大学 ｜ 电子科学与技术（本科）}
\hfill \daterange{2021.09 \textasciitilde{} 2025.06}

\par\noindent
$\bullet$ \textbf{GPA：}3.61 / 5.0（专业前 20\%） \hspace{0.5em} \textbf{主修课程：}数据结构、计算机组成原理、微机原理、半导体物理与器件等。

\section*{实习经历}
\vspace{-0.1em}
\noindent\entrytitle{北京开源芯片研究院 ｜ SoC 项目组嵌入式软件实习生}\hfill\daterange{2025.02 \textasciitilde{} 2025.07}
\par\noindent $\bullet$ \textbf{CPU 测试与工具链：}开展项目组 CPU 核功能调研与测试，研究代码压缩方案，参与编译器工具链调优。
\par\noindent $\bullet$ \textbf{车规 ECU 虚拟化：}围绕 ARM Cortex-R52 与 Bao Hypervisor，分析特权级、内存保护、中断注入与 vCPU 调度机制，梳理 RISC-V 移植所需的架构能力与功能取舍。
\section*{科研经历}
\vspace{-0.1em}
\noindent\entrytitle{NACRE ｜ 机密容器的原生体系结构支持}\hfill\daterange{2026.08 \textasciitilde{} 至今}
\par\noindent $\bullet$ \textbf{架构机制与硬件移植：}将 QEMU 原型中的相关机制移植至 BOOM 处理器，结合硬件实现需求调整设计；基于 FireSim 完成 FPGA 原型部署。
\par\addvspace{0.45\baselineskip}
\noindent\entrytitle{CPU Fuzzing ｜ RISC-V ISA 层错误检测}\hfill\daterange{2025.09 \textasciitilde{} 至今}
\par\noindent $\bullet$ \textbf{架构正确性研究：}面向 RISC-V 处理器开展模糊测试，依据 ISA 规范检测处理器架构行为与规范之间的偏差。
\section*{项目经历}
\vspace{-0.1em}
\noindent\entrytitle{“一生一芯”（ysyx）｜ RISC-V CPU/SoC 软硬件系统}\hfill\daterange{2023.09 \textasciitilde{} 2025.09}
\par\noindent $\bullet$ \textbf{RTL 实现与差分验证：}基于 Chisel 实现 RISC-V CPU，使用 Verilator 进行 RTL 仿真，通过自研 Difftest 与指令集模拟器开展差分验证。
\par\noindent $\bullet$ \textbf{ISA 模拟器与架构功能：}实现 RV32IM 指令集模拟器，以 Spike 为参考验证指令行为；扩展 CSR 与页表支持，为系统软件运行提供架构功能。
\par\noindent $\bullet$ \textbf{系统软件运行：}移植并启动 RT-Thread；在自研模拟器上通过 OpenSBI 启动 Linux，进入 Shell 交互环境。
\par\addvspace{0.45\baselineskip}
\noindent\entrytitle{集成电路创新创业大赛 ｜ 存储控制器与图像压缩 IP}\hfill\daterange{2022.03 \textasciitilde{} 2024.06}
\par\noindent $\bullet$ \textbf{多端口存储控制器：}完成 16×16 多端口 SRAM 控制器的 Verilog RTL 实现与 SVA 验证；45nm 工艺下主频 250MHz，总带宽 16Gbps。
\par\noindent $\bullet$ \textbf{JPEG-LS 与 SoC 集成：}参与 C/C++ 模型、Verilog IP 实现与 AXI4 集成，在 MPSoC 平台完成 FPGA 回环测试；获 2023 年省级二等奖。
\section*{专业技能与荣誉}
\vspace{-0.1em}
\par\noindent $\bullet$ \textbf{处理器设计与验证：}Chisel/Verilog、Verilator、SVA、FPGA 原型验证；ISA 模拟器、Spike、Difftest 与 CPU Fuzzing。
\par\noindent $\bullet$ \textbf{体系结构与系统软件：}RISC-V、CSR、页表；RT-Thread 移植、OpenSBI/Linux 启动；QEMU、BOOM 与 FireSim 原型部署。
\par\noindent $\bullet$ \textbf{编程与开发环境：}C/C++、Python、Scala；Linux、Git、Makefile。
\par\noindent $\bullet$ \textbf{荣誉奖项：}EDA 精英挑战赛国家一等奖、电子设计大赛省一等奖、集成电路创新创业大赛省二等奖；二等奖学金、优秀学生干部。
\section*{校园经历}
\vspace{-0.1em}
\noindent\entrytitle{南京邮电大学大学生科协电子部 ｜ 部长 / 讲师}\hfill\daterange{2022.09 \textasciitilde{} 2023.06}
\par\noindent $\bullet$ 负责社团日常运营与新生授课，策划校“科技节”比赛并保障赛程推进。

\end{document}
